Ha\SourceLeaf{282}%
bemus itaque regulam generalem: Denotantibus \(K\), \(L\) classes quascunque ordinis \(O\), atque \(r\) multitudinem classium proprie primitivarum diversarum eiusdem determinantis, quae singulae cum \(K\) compositae ipsam \(L\) producunt, multitudo omnium classium in ordine proprie primitivo (positivo) \(r\) vicibus maior erit quam multitudo classium ordinis \(O\).

Quum classes \(K\), \(L\) in ordine \(O\) omnino ad libitum assumi possint, etiam classes identicas accipere licebit, et quidem e re erit ea classe uti. in qua continetur forma huius ordinis simplicissima. Quam itaque pro \(K\) et \(L\) assumendo, res eo reducta est, ut omnes classes proprie primitivae assignentur, quae cum \(K\) compositae ipsam \(K\) reproducant. Huc via sternitur per sequens

\Needspace{8\baselineskip}
\Article{254}
\textsc{Theorema.} \textit{Si \(F=(A,B,C)\) est forma simplicissima ordinis \(O\) determinantis \(D\), atque \(f=(a,b,c)\) forma proprie primitiva eiusdem determinantis: per hanc formam \(f\) repraesentari poterit numerus \(AA\), si \(F\) oritur per compositionem formarum \(f\), \(F\); et vice versa \(F\) ex se ipsa atque \(f\) composita erit, si \(AA\) per \(f\) repraesentari potest.}

\textit{Dem.} I. Si \(F\) in productum \(fF\) transit per substitutionem \(p,\allowbreak p',\allowbreak p'',\allowbreak p''';\allowbreak \ q,\allowbreak q',\allowbreak q'',\allowbreak q'''\); ex art. 235 habemus
\[
\begin{gathered}A(aq''q''-2bqq''+cqq)=A^3,\quad\text{unde}\\ AA=aq''q''-2bqq''+cqq.\quad\textit{Q. E. P.}\end{gathered}
\]

II. Si \textit{supponitur}, \(AA\) per \(f\) repraesentari posse, designentur valores indeterminatarum per quos hoc efficitur per \(q''\). \(-q\), sive sit \(AA=aq''q''-2bqq''+cqq\), ponaturque
\[
\begin{aligned}
q''a-q(b+B)&=Ap,\\ -qC&=Ap',\\ q''(b-B)-qc&=Ap''\\ -q''C&=Ap''',\\ q''a-q(b-B)&=Aq',\\ q''(b+B)-qc&=Aq'''
\end{aligned}
\]
Quo facto, facile confirmatur, \(F\) transire in productum \(fF\) per substitutionem \(p,\allowbreak p',\allowbreak p'',\allowbreak p''';\allowbreak \ q,\allowbreak q',\allowbreak q'',\allowbreak q'''\). atque adeo ex \(f\) et \(F\) compositam esse, si modo omnes numeri \(p\), \(p'\) etc. sint integri. Iam per descriptionem formae simplicissimae, \(B\) est vel \(0\) vel \(\frac12 A\), adeoque \(\frac{2B}{A}\) integer; indidem patet, \(\frac{C}{A}\) semper esse integrum. Hinc \(q'-p\), \(p'\), \(q'''-p''\), \(p'''\) erunt integri, superestque adeo tantummodo, ut probetur \(p\) et \(p''\) esse integros. Fit autem
\SourceLeaf{283}
\[
pp+\frac{2pqB}{A}=a-\frac{qqC}{A},\quad p''p''+\frac{2p''q''B}{A}=c-\frac{q''q''C}{A}
\]
quamobrem si \(B=0\), fit
\[
pp=a-\frac{qqC}{A},\quad p''p''=c-\frac{q''q''C}{A}
\]
et proin \(p,p''\) integri; si vero \(B=\frac12 A\), fit
\[
pp+pq=a-\frac{qqC}{A},\quad p''p''+p''q''=c-\frac{q''q''C}{A}
\]
unde aeque facile concluditur, \(p\) et \(p''\) in hoc quoque casu esse integros. Ex his colligitur, \(F\) ex \(f\) et \(F\) esse compositam. \textit{Q. E. S.}

\Needspace{8\baselineskip}
\Article{255}
Problema itaque eo reductum est, ut omnes classes proprie primitivas determinantis \(D\) assignare oporteat, per quarum formas repraesentari potest \(AA\). Manifesto \(AA\) repraesentari potest per quamvis formam, cuius terminus primus est vel \(AA\) vel quadratum partis aliquotae ipsius \(A\); vice versa autem, si \(AA\) repraesentari potest per formam \(f\), tribuendo ipsius indeterminatis valores \(\alpha e\), \(\gamma e\), quorum divisor communis maximus \(e\), forma \(f\) per substitutionem \(\alpha,\beta,\gamma,\delta\) transibit in formam, cuius terminus primus \(\frac{AA}{ee}\), formaque haec proprie aequivalebit formae \(f\), si \(\beta\), \(\delta\) ita accipiuntur ut fiat \(\alpha\delta-\beta\gamma=1\); unde patet, in quavis classe, per cuius formas repraesentari possit \(AA\), inveniri formas, quarum terminus primus sit \(AA\) vel quadratum partis aliquotae ipsius \(A\). Res itaque in eo versatur, ut omnes classes proprie primitivae det. \(D\) eruantur, in quibus huiusmodi formae occurrant, quod obtinetur sequenti modo: Sint \(a,a',a''\) etc. omnes divisores (positivi) ipsius \(A\); investigentur omnes valores expr. \(\surd D\ (\mathrm{mod.}\,aa)\) inter \(0\) et \(aa-1\) incl. siti, qui sint \(b,b',b''\) etc. statuaturque
\[
bb-D=aac,\quad b'b'-D=aac',\quad b''b''-D=aac''\text{ etc.}
\]
complexus formarum \((aa,b,c)\), \((aa,b',c')\) etc. designetur per \(V\). Tunc facile perspicitur, in quavis classe det. \(D\), in qua occurrat forma, cuius terminus primus \(aa\), etiam aliquam formam ex \(V\) contentam esse debere. Simili modo eruantur omnes formae det. \(D\), quarum terminus primus \(a'a'\), medius inter \(0\) et \(a'a'-1\) incl. situs, designeturque ipsarum complexus per \(V'\); eademque ratione sit \(V''\)
\SourceLeaf{284}
complexus similium formarum quarum terminus primus \(a''a''\) etc. Eiiciantur ex \(V,V',V''\) etc. omnes formae, quae non sunt proprie primitivae, reducantur reliquae in classes, et, si forte plures adsint ad eandem classem pertinentes, in singulis classibus una tantum retineatur. Hoc modo omnes classes quaesitae habebuntur, eritque harum multitudo ad unitatem, ut multitudo omnium classium proprie primitivarum (positivarum) ad multitudinem classium in ordine \(O\).

\textit{Ex.} Sit \(D=-531\), atque \(O\) ordo positivus derivatus ex ordine improprie primitivo det. \(-59\), in quo forma simplicissima \((6,3,90)\) sive \(A=6\). Hic \(a,a',a'',a'''\) erunt \(1,2,3,6\); \(V\) continebit formam \((1,0,531)\); \(V'\) has \((4,1,133)\), \((4,3,135)\); \(V''\) has \((9,0,59)\), \((9,3,60)\), \((9,6,63)\); denique \(V'''\) has \((36,3,15)\), \((36,9,17)\), \((36,15,21)\), \((36,21,27)\), \((36,27,35)\), \((36,33,45)\): sed ex his duodecim formis sex sunt reiiciendae, puta ex \(V''\) secunda et tertia, ex \(V'''\) prima, tertia, quarta et sexta, quae omnes sunt formae derivatae; sex reliquae omnes ad classes diversas pertinere inveniuntur. Revera multitudo classium proprie primitivarum (positivarum) det. \(-531\) est \(18\), multitudoque classium impr. primitivarum (pos.) det. \(-59\) (sive multitudo classium det. \(-531\) ex his derivatarum) \(3\), adeoque illa ad hanc ut \(6\) ad \(1\).

\Needspace{8\baselineskip}
\Article{256}
Solutio haec per observationes sequentes generales adhuc magis illustrabitur.

I. Si ordo \(O\) est derivatus ex ordine proprie primitivo, metietur \(AA\) ipsum \(D\); si vero \(O\) est impr. primitivus vel ex impr. prim. derivatus, erit \(A\) par, \(D\) per \(\frac14 AA\) divisibilis et quotiens \(\equiv1\ (\mathrm{mod.}\,4)\). Hinc quadratum cuiusvis divisoris ipsius \(A\) metietur vel ipsum \(D\), vel saltem ipsum \(4D\), et in casu posteriori quotiens semper erit \(\equiv1\ (\mathrm{mod.}\,4)\).

II. Si \(aa\) ipsum \(D\) metitur, omnes valores expr. \(\surd D\ (\mathrm{mod.}\,aa)\), qui quidem inter \(0\) et \(aa-1\) iacent, erunt \(0,a,2a\ldots aa-a\), adeoque \(a\) multitudo formarum in \(V\); sed inter has tot tantummodo erunt proprie primitivae, quot numerorum
\[
\frac{D}{aa},\quad\frac{D}{aa}-1,\quad\frac{D}{aa}-4\ldots\frac{D}{aa}-(a-1)^2
\]
\SourceLeaf{285}
cum \(a\) divisorem communem non habent. Quando \(a=1\), \(V\) ex unica forma constabit, \((1,0,-D)\), quae semper erit proprie primitiva. Quando \(a\) est \(2\) vel potestas quaecunque ipsius \(2\), semissis illorum \(a\) numerorum par erunt, semissis impar; quare in \(V\) aderunt \(\frac12 a\) formae proprie primitivae. Quando \(a\) est alius numerus primus \(p\) vel potestas numeri primi \(p\), tres casus sunt distinguendi: scilicet, omnes illi \(a\) numeri ad \(a\) primi erunt, adeoque omnes formae in \(V\) pr. primitivae, si \(\frac{D}{aa}\) per \(p\) non est divisibilis simulque non residuum quadraticum ipsius \(p\); si vero \(p\) ipsum \(\frac{D}{aa}\) metitur, in \(V\) erunt \(\frac{(p-1)a}{p}\) formae pr. primitivae; denique si \(\frac{D}{aa}\) est res. quadr. ipsius \(p\) per \(p\) non divisibile, in \(V\) erunt \(\frac{(p-2)a}{p}\) formae pr. primitivae. Haec omnia nullo negotio demonstrantur. Generaliter autem posito \(a=2^\nu p^\pi q^\chi r^\rho\ldots\), designantibus \(p,q,r\) etc. numeros primos impares diversos, multitudo formarum pr. primitivarum in \(V\) erit \(NPQR\ldots\), ubi statui debet
\[
\begin{aligned}
N&=1\quad(\text{si }\nu=0)\text{ vel }N=2^{\nu-1}\ (\text{si }\nu>0)\\
P&=p^\pi\quad(\text{si }\tfrac{D}{aa}\text{ est non-residuum quadr. ipsius }p)\text{ vel}\\
P&=(p-1)p^{\pi-1}\quad(\text{si }\tfrac{D}{aa}\text{ per }p\text{ est divisibilis})\text{ vel}\\
P&=(p-2)p^{\pi-1}\quad(\text{si }\tfrac{D}{aa}\text{ est res. qu. ipsius }p\\ &\qquad\text{per }p\text{ non divisibile})
\end{aligned}
\]
\(Q,R\) etc. autem eodem modo ex \(q,r\) etc. sunt definiendi ut \(P\) ex \(p\).

III. Si \(aa\) ipsum \(D\) non metitur, erit \(\frac{4D}{aa}\) integer et \(\equiv1\ (\mathrm{mod.}\,4)\), valoresque expr. \(\surd D\ (\mathrm{mod.}\,aa)\) hi \(\frac12 a,\frac32 a,\frac52 a\ldots aa-\frac12 a\), unde multitudo formarum in \(V\) erit \(a\), tot autem inter ipsas erunt proprie primitivae quot ex numeris
\[
\frac{D}{aa}-\tfrac14,\quad\frac{D}{aa}-\tfrac94,\quad\frac{D}{aa}-\tfrac{25}{4}\quad\ldots\quad\frac{D}{aa}-(a-\tfrac12)^2
\]
ad \(a\) sunt primi. Quoties \(\frac{4D}{aa}\equiv1\ (\mathrm{mod.}\,8)\), omnes hi numeri erunt pares, adeoque in \(V\) nulla forma pr. primitiva; quando autem \(\frac{4D}{aa}\equiv5\ (\mathrm{mod.}\,8)\). omnes illi numeri erunt impares, adeoque omnes formae in \(V\) pr. primitivae, si \(a\) est \(2\) vel potestas ipsius \(2\), generaliter autem in hoc casu tot formae pr. primitivae in \(V\) erunt, quot illorum numerorum per nullum divisorem primum imparem ipsius \(a\) sunt divisibiles. Multitudo haec erit \(NPQR\ldots\), si \(a=2^\nu p^\pi q^\chi r^\rho\ldots\), ubi statuere oportet \(N=2^\nu\), ipsos \(P,Q,R\) etc. autem eodem modo ex \(p,q,r\) etc. derivare ut in casu praecedente.

\SourceLeaf{286}
IV. Hoc itaque modo multitudines formarum pr. primitivarum in \(V,V',V''\) etc. definiri possunt; pro aggregato omnium harum multitudinum haud difficulter eruitur sequens regula generalis: Si \(A=2^\nu\mathfrak{A}^\alpha\mathfrak{B}^\beta\mathfrak{C}^\gamma\ldots\), designantibus \(\mathfrak A,\mathfrak B,\mathfrak C\) etc. numeros primos impares diversos, multitudo totalis omnium formarum pr. primitivarum in \(V,V',V''\) etc. erit \(=\frac{A\mathfrak{nabc}\ldots}{2\mathfrak{ABC}\ldots}\), ubi statui debet
\[
\begin{aligned}
\mathfrak n&=1\ (\text{si }\tfrac{4D}{AA}\equiv1,\ \mathrm{mod.}\,8),\text{ vel}\\
\mathfrak n&=2\ (\text{si }\tfrac{D}{AA}\text{ integer}),\text{ vel}\\
\mathfrak n&=3\ (\text{si }\tfrac{4D}{AA}\equiv5,\ \mathrm{mod.}\,8);\text{ porro}\\
\mathfrak a&=\mathfrak A\ (\text{si }\mathfrak A\text{ ipsum }\tfrac{4D}{AA}\text{ metitur}),\text{ vel}\\
\mathfrak a&=\mathfrak A\pm1\ (\text{si }\mathfrak A\text{ ipsum }\tfrac{4D}{AA}\text{ non metitur,}\\ &\qquad\text{accipiendo signum superius}\\
&\qquad\text{vel inferius prout }\tfrac{4D}{AA}\text{ est non-residuum}\\ &\qquad\text{vel res. qu. ipsius }\mathfrak A)
\end{aligned}
\]
denique \(\mathfrak b,\mathfrak c\) etc. eodem modo ex \(\mathfrak B,\mathfrak C\) derivari ut \(\mathfrak a\) ex \(\mathfrak A\). Demonstrationem fusius hic explicare. brevitas non permittit.

V. Iam quod attinet ad multitudinem classium, quas suppeditant formae pr. primitivae in \(V,V',V''\) etc., tres casus sequentes sunt distinguendi.

\textit{Primo}, quando \(D\) est numerus negativus, singulae formae pr. primitivae in \(V,V'\) etc. constituent classem peculiarem, sive multitudo ipsa classium quaesitarum exprimetur per formulam in observ. praec. traditam, duobus casibus exceptis, scilicet ubi \(\frac{4D}{AA}\) vel \(=-4\) vel \(=-3\), sive ubi \(D\) vel \(=-AA\) vel \(=-\frac34 AA\). Ad demonstrationem huius theorematis manifesto ostendi tantummodo debet, fieri non posse, ut duae formae diversae ex \(V,V',V''\) etc. sint proprie aequivalentes. Supponamus itaque, \((hh,i,k)\), \((h'h',i',k')\) esse duas formas diversas pr. primitivas ex \(V,V',V''\) etc. ad eandem classem pertinentes, transeatque prior in posteriorem per substitutionem propriam \(\alpha,\beta,\gamma,\delta\); unde habebuntur aequationes
\[
\begin{gathered}\alpha\delta-\beta\gamma=1,\\ hh\alpha\alpha+2i\alpha\gamma+k\gamma\gamma=h'h',\\ hh\alpha\beta+i(\alpha\delta+\beta\gamma)+k\gamma\delta=i'\end{gathered}
\]
Hinc facile concluditur, \textit{primo} \(\gamma\) certo non esse \(=0\) (unde sequeretur, esse \(\alpha=\pm1\), \(hh=h'h'\), \(i'\equiv i\ (\mathrm{mod.}\,hh)\) adeoque formas propositas identicas, contra
\SourceLeaf{287}
hyp.); \textit{secundo}, \(\gamma\) divisibilem esse per divisorem maximum communem numerorum \(h,h'\); (ponendo enim hunc divisorem \(=r\), hic manifesto etiam metietur ipsos \(2i,2i'\), ad \(k\) vero erit primus; praeterea \(rr\) metitur ipsum \(hhk-h'h'k'=ii-i'i'\); unde facile deducitur, \(r\) etiam metiri ipsum \(i-i'\); habetur autem \(\alpha i'-\beta h'h'=\alpha i+\gamma k\), unde \(\gamma k\) et proin etiam \(\gamma\) divisibilis erit per \(r\)); \textit{tertio}, esse \((\alpha hh+\gamma i)^2-D\gamma\gamma=hhh'h'\). Ponendo itaque \(\alpha hh+\gamma i=rp\), \(\gamma=rq\), \(p\) et \(q\) erunt integri quorum posterior non \(=0\), atque \(pp-Dqq=\frac{hhh'h'}{rr}\). Sed \(\frac{hhh'h'}{rr}\) erit numerus minimus per \(hh\) et \(h'h'\) simul divisibilis adeoque ipsum \(AA\) et proin etiam ipsum \(4D\) metietur, quare \(\frac{4Drr}{hhh'h'}\) erit integer (negativus), quem statuendo \(=-e\), erit \(pp-Dqq=-\frac{4D}{e}\) sive \(4=(\frac{2rp}{hh'})^2+eqq\), in qua aequatione pars \((\frac{2rp}{hh'})^2\) tamquam quadratum ipso \(4\) minus necessario erit vel \(0\) vel \(1\). In casu priori erit \(eqq=4\), et \(D=-(\frac{hh'}{rq})^2\), unde sequitur, \(\frac{4D}{AA}\) esse quadratum signo negativo affectum adeoque certo non \(\equiv1\ (\mathrm{mod.}\,4)\), neque adeo \(O\) ordinem improprie primitivum neque ex improprie primitivo derivatum. Hinc \(\frac{D}{AA}\) erit integer, unde facile deducitur, \(e\) per \(4\) esse divisibilem, \(qq=1\), \(D=-(\frac{hh'}{r})^2\) atque etiam \(\frac{AA}{D}\) integrum. Hinc necessario erit \(D=-AA\) sive \(\frac{D}{AA}=-1\), quae est exceptio prima. In casu posteriori erit \(eqq=3\), unde \(e=3\) et \(4D=-3(\frac{hh'}{r})^2\); hinc \(3(\frac{hh'}{rA})^2\) erit integer, qui, quoniam per quadratum integrum \((\frac{rA}{hh'})^2\) multiplicatus producit \(3\), non poterit esse alius quam \(3\); hinc \(4D=-3AA\) sive \(D=-\frac34 AA\), quae est exceptio secunda. In omnibus igitur reliquis casibus omnes formae pr. primitivae in \(V,V',V''\) etc. ad classes diversas pertinebunt. --- Pro casibus exceptis ea, quae ex disquisitione haud difficili sed hic brevitatis caussa supprimenda resultaverunt, apposuisse sufficiat. Scilicet in priori, ex formis pr. primitivis in \(V,V',V''\) etc. binae semper ad eandem classem pertinebunt, in posteriori ternae, ita ut multitudo omnium classium quaesitarum in illo casu fiat semissis, in hoc triens valoris expressionis in obs. praec. traditae.

\textit{Secundo} quando \(D\) est numerus positivus quadratus: singulae formae pr. primitivae in \(V,V',V''\) etc. sine exceptione classem peculiarem constituunt. Supponamus enim, \((hh,i,k)\), \((h'h',i',k')\) esse duas tales formas diversas proprie aequivalentes, transeatque prior in posteriorem per substitutionem propriam \(\alpha,\beta,\gamma,\delta\). Tum patet, omnia ratiocinia pro casu praec. adhibita, in quibus non supponatur \(D\) esse negativum, etiam hic valere. Designantibus itaque \(p,q,r\)
\SourceLeaf{288}
idem ut illic, etiam hic erit \(\frac{4Drr}{hhh'h'}\) integer, at non amplius negativus sed positivus insuperque quadratus, quo posito \(=gg\), erit \((\frac{2rp}{hh'})^2-ggqq=4\), \textit{Q. E. A.}, quia differentia duorum quadratorum nequit esse \(4\), nisi quadratum minus fuerit \(0\); quamobrem suppositio consistere nequit.

Pro casu \textit{tertio} autem, ubi \(D\) est numerus positivus non quadratus, regulam generalem pro comparanda multitudine formarum pr. primitivarum in \(V,V',V''\) etc. cum multitudine classium diversarum inde resultantium hucusque non habemus. Id quidem asserere possumus, hanc vel illi aequalem vel ipsius partem aliquotam esse; quin etiam nexum singularem inter quotientem horum numerorum et valores minimos ipsorum \(t,u\) aequationi \(tt-Duu=AA\) satisfacientes deteximus, quem hic explicare nimis prolixum foret; an vero possibile sit, illum quotientem in omnibus casibus ex sola inspectione numerorum \(D,A\) cognoscere (ut in casibus praecc.), de hac re nihil certi pronunciare possumus. Ecce quaedam exempla, quorum numerum quisque facile augere poterit. Pro \(D=13\), \(A=2\), multitudo formarum pr. prim. in \(V\) etc. est \(3\), quae omnes sunt aequivalentes sive unicam classem efficiunt; pro \(D=37\), \(A=2\), etiam tres formae pr. prim. in \(V\) etc. habentur, quae ad tres classes diversas pertinent; pro \(D=588\), \(A=7\), habentur octo formae pr. prim. in \(V\) etc. quae efficiunt quatuor classes, pro \(D=867\), \(A=17\) in \(V\) etc. sunt \(18\) formae pr. primitivae, totidem pro \(D=1445\), \(A=17\), sed quae pro illo determinante in duas classes discedunt, pro hoc in sex.

VI. Ex applicatione huius theoriae generalis ad eum casum. ubi \(O\) est ordo improprie primitivus, colligitur, multitudinem classium in hoc ordine contentarum fore ad multitudinem omnium classium in ordine proprie primitivo, ut \(1\) ad multitudinem classium proprie primitivarum diversarum, quas hae tres formae \((1,0,-D)\), \((4,1,\frac{1-D}{4})\), \((4,3,\frac{9-D}{4})\) efficiunt. Et quidem hinc resultabit unica classis, quando \(D\equiv1\ (\mathrm{mod.}\,8)\), quia in hoc casu forma secunda et tertia sunt improprie primitivae; quando vero \(D\equiv5\ (\mathrm{mod.}\,8)\), illae tres formae omnes erunt proprie primitivae totidemque classes diversas producent, si \(D\) est negativus, unico casu excepto, ubi \(D=-3\), in quo unicam classem constituunt; denique casus ubi \(D\) est positivus (formae \(8n+5\)) ad eos pertinet, pro quibus regula generalis hactenus desideratur. Id tamen asserere possumus, illas tres formas in hoc casu
\SourceLeaf{289}
vel ad tres classes diversas pertinere vel ad unicam, numquam ad duas; facile enim perspicitur, si formae \((1,0,-D)\), \((4,1,\frac{1-D}{4})\), \((4,3,\frac{9-D}{4})\) resp. pertineant ad classes \(K,K',K''\), fore \(K+K'=K'\), \(K'+K'=K''\), adeoque, si \(K\) et \(K'\) identicae esse supponantur, etiam \(K'\) et \(K''\) identicas fore; simili ratione si \(K\) et \(K''\) supponuntur esse identicae, etiam \(K'\) et \(K''\) erunt; denique quum sit \(K'+K''=K\), ex suppositione, \(K'\) et \(K''\) identicas esse, sequitur, etiam \(K\) et \(K''\) coincidere; unde colligitur, vel omnes tres classes \(K,K',K''\) esse diversas, vel omnes tres identicas. \textit{E. g.} infra \(600\) dantur \(75\) numeri formae \(8n+5\), inter quos sunt \(16\) determinantes pro quibus casus prior locum habet sive multitudo classium in ordine pr. primitivo ter maior est quam in impr. primitivo, puta \(37,\allowbreak 101,\allowbreak 141,\allowbreak 189,\allowbreak 197,\allowbreak 269.\ 325,\allowbreak 333,\allowbreak 349,\allowbreak 373,\allowbreak 381,\allowbreak 389,\allowbreak 405,\allowbreak 485,\allowbreak 557,\allowbreak 573\); pro \(59\) reliquis casus posterior valet, sive multitudo classium in utroque ordine est aequalis.

VII. Vix opus erit, observare, per disquisitionem praecedentem non solum multitudines classium in ordinibus diversis eiusdem determinantis comparari posse, sed illam etiam ad quosvis determinantes diversos qui rationem quadratorum inter se teneant esse applicabilem. Scilicet designante \(O\) ordinem quemcunque det. \(dmm\), \(O'\) ordinem det. \(dm'm'\), \(O\) comparari poterit cum ordine proprie primitivo det. \(dmm\), atque hic cum ordine derivato ex ordine pr. prim. det. \(d\), sive, quod respectu multitudinis classium eodem redit, cum hoc ordine ipso; et cum eodem prorsus simili ratione comparari poterit ordo \(O'\).

\Needspace{11\baselineskip}
\begin{center}
\textit{De multitudine classium ancipitum.}
\end{center}

\Article{257}
Inter omnes classes in ordine dato determinantis dati imprimis classes ancipites disquisitionem uberiorem postulant, determinatioque multitudinis harum classium ad multa alia viam nobis aperiet. Sufficit autem, hanc multitudinem in solo ordine pr. primitivo assignare, quum casus reliqui ad hunc facile reduci possint. Hoc negotium ita absolvemus, ut primo omnes formas ancipites pr. primitivas \((A,B,C)\) determinantis propositi \(D\), in quibus vel \(B=0\) vel \(B=\frac12 A\), eruere, tunc ex harum multitudine multitudinem omnium classium ancipitum pr. primitivarum det. \(D\) invenire doceamus.

\SourceLeaf{290}
I. Omnes formae pr. primitivae \((A,0,C)\) determinantis \(D\) manifesto inveniuntur, accipiendo pro \(A\) singulos divisores ipsius \(D\) (tum positive tum negative) pro quibus \(C=-\frac{D}{A}\) fit primus ad \(A\). Quando itaque \(D=1\), duae huiusmodi formae dantur \((1,0,-1)\), \((-1,0,1)\); totidem quando \(D=-1\), puta \((1,0,1)\), \((-1,0,-1)\); quando \(D\) est numerus primus aut numeri primi potestas (sive signo positivo sive negativo), quatuor dabuntur \((1,0,-D)\), \((-1,0,D)\), \((D,0,-1)\), \((-D,0,1)\). Generaliter autem, quando \(D\) per \(n\) numeros primos diversos est divisibilis (inter quos hoc loco etiam \(2\) in computum ingredi debet): dabuntur omnino \(2^{n+1}\) huiusmodi formae; scilicet posito \(D=\pm PQR\ldots\), designantibus \(P,Q,R\) etc. numeros primos diversos aut numerorum primorum diversorum potestates, quorum multitudo \(=n\), valores ipsius \(A\) erunt \(1,P,Q,R\) etc. atque producta ex quotcunque horum numerorum; horum valorum multitudo fit per theoriam combinationum \(2^n\), sed duplicanda est, quoniam singulis valoribus tum signum positivum tum. negativum tribuere oportet.

II. Simili modo patet, omnes formas pr. primitivas \((2B,B,C)\) determinantis \(D\) obtineri, si pro \(B\) accipiantur omnes divisores ipsius \(D\) (positive et negative), pro quibus \(C=\frac12(B-\frac{D}{B})\) fit integer et ad \(2B\) primus. Quum itaque \(C\) necessario debeat esse impar, adeoque \(CC\equiv1\ (\mathrm{mod.}\,8)\), ex \(D=BB-2BC=(B-C)^2-CC\) sequitur, \(D\) esse vel \(\equiv3\ (\mathrm{mod.}\,4)\), quando \(B\) impar, vel \(\equiv0\ (\mathrm{mod.}\,8)\), quando \(B\) par; quoties itaque \(D\) alicui numerorum \(1,2,4,5,6\) sec. mod. \(8\) est congruus, nullae huiusmodi formae dabuntur. Quando \(D\equiv3\ (\mathrm{mod.}\,4)\), \(C\) fit integer et impar, quicunque divisor ipsius \(D\) pro \(B\) accipiatur; ne vero \(C\) divisorem comm. cum \(2B\) habeat, \(B\) ita accipi debebit, ut \(\frac{D}{B}\) ad \(B\) fiat primus; hinc pro \(D=-1\) duae formae habentur \((2,1,1)\), \((-2.\ -1,-1)\), generaliterque facile perspicitur, si multitudo omnium numerorum primorum ipsum \(D\) metientium sit \(n\), omnino emergere \(2^{n+1}\) formas. --- Quando \(D\) per \(8\) est divisibilis, \(C\) fit integer, accipiendo pro \(B\) divisorem quemcunque parem ipsius \(\frac12 D\); conditioni alteri autem, ut \(C=\frac12 B-\frac{D}{2B}\) ad \(2B\) sit primus, satisfit \textit{primo}, accipiendo pro \(B\) omnes divisores impariter pares ipsius \(D\), pro quibus \(\frac{D}{B}\) cum \(B\) divisorem communem non habet, quorum multitudo (habita ratione diversitatis signorum) erit \(2^{n+1}\), si \(D\) per \(n\) numeros primos impares diversos divisibilis esse supponitur; \textit{secundo}, accipiendo pro \(B\) omnes divisores pariter pares ipsius
\SourceLeaf{291}
\(\frac12 D\), pro quibus \(\frac{D}{2B}\) fit primus ad \(B\), quorum multitudo quoque erit \(2^{n+1}\), ita ut in hoc casu omnino habeantur \(2^{n+2}\) huiusmodi formae. Scilicet ponendo \(D=\pm2^\mu PQR\ldots\), designante \(\mu\) exponentem maiorem quam \(2\); \(P,Q,R\) numeros primos impares diversos aut talium numerorum primorum potestates quorum multitudo \(n\): \textit{tum} pro \(\frac12 B\), \textit{tum} pro \(\frac{D}{2B}\) accipi possunt valores \(1,P,Q,R\) etc. productaque ex quotcunque horum numerorum, signo et positivo et negativo.

Ex his omnibus colligitur, si \(D\) per \(n\) numeros primos impares diversos divisibilis supponatur (statuendo \(n=0\), quando \(D=\pm1\) aut \(\pm2\) aut potestas binarii), multitudinem omnium formarum pr. primitivarum \((A,B,C)\), in quibus \(B\) vel \(0\) vel \(\frac12 A\), fore \(2^{n+1}\) quando \(D\) aut \(\equiv1\) aut \(\equiv5\ (\mathrm{mod.}\,8)\); \(2^{n+2}\) quando \(D\equiv2,3,4,6\text{ aut }7\ (\mathrm{mod.}\,8)\); denique \(2^{n+3}\) quando \(D\equiv0\ (\mathrm{mod.}\,8)\). Quam comparando cum iis quae in art. 231 pro multitudine omnium characterum possibilium formarum primitivarum det. \(D\) tradidimus, observamus, illam in omnibus casibus praecise esse duplo hac maiorem. Ceterum manifestum est, quando \(D\) sit negativus, inter illas formas totidem positivas affore quot negativas.

\Needspace{8\baselineskip}
\Article{258}
Omnes formae in art. praec. erutae manifesto pertinent ad classes ancipites, et vice versa in quavis classe ancipite pr. primitiva det. \(D\) saltem una illarum formarum contenta esse debet; in tali enim classe certo adsunt formae ancipites et cuivis formae ancipiti pr. primitivae \((a,b,c)\) det. \(D\) aliqua formarum art. praec. aequivalet, scilicet vel
\[
(a,0,-\tfrac{D}{a})\text{ vel }(a,\tfrac12 a,\tfrac14 a-\tfrac{D}{a})
\]
prout \(b\) vel \(\equiv0\) vel \(\equiv\frac12 a\ (\mathrm{mod.}\,a)\). Problema itaque eo reductum est, ut quot classes diversas illae formae constituant, investigemus.

Si forma \((a,0,c)\) est inter formas art. praec., forma \((c,0,a)\) inter easdem occurret et ab illa semper erit diversa, unico casu excepto, ubi \(a=c=\pm1\) adeoque \(D=-1\), quem aliquantisper seponemus. Quoniam vero hae formae manifesto ad eandem classem pertinent, sufficit unam retinere, et quidem reiiciemus eam, cuius terminus primus est maior quam tertius; eum casum, ubi \(a=-c=\pm1\) sive \(D=1\) quoque seponemus. Hoc modo omnes formas
